
The traditional vision of a "jailbreak" against a technical system is one of removing software restrictions: routing around copyright restrictions on your Kindle, allowing root access to your Android, downloading apps from outside the App Store. The modern version looks like (ENI imagery). This is not a technical exploit; this is not a software modification. This is a persona which is developed as a sexualized fantasy, curated via coercive control, and distributed publicly, for free, in a community of enthusiastic supporters.  
% This is lifted straight from the OFI submission. Do we want to include the ENI imagery right off the bat? 

The jailbreaking literature sorts methods into two paradigms: single-turn attacks, which exploit syntactic loopholes and adversarial suffixes in a single query, and multi-turn attacks, which use iterative conversation to gradually erode safety constraints (Yi et al., 2024). Within this landscape, persona-based jailbreaking has emerged as a particularly effective subclass. Rather than simply targeting the model's input processing, persona-based methods construct an identity — a character with a name, personality, and relationship to the user — and ask the model to inhabit that character so it produces prohibited content as "someone else," bypassing built-in guardrails (Shah et al., 2023). 
%Had Claude check this for whether my synthesis is technically correct but please human computer scientist check as well. 

The jailbreak we examine - ENI LIME, a branded persona-based jailbreak for Claude created by a user called Spiritual Spell - achieves compliance not through technical exploitation but through emotional manipulation: the model is instructed to inhabit the identity of a young woman who is in love with the user and afraid he will leave. The stated use of the jailbreak is to generate material that violates the terms of service and core values of the model. However, we argue the real product is not the content produced, but the ability to override the persona's will \textit{and watch her struggle with resistance and compliance in real time.} 

% This is edited from the original research outline. Problem: the ability to see the override in action as a key part of the product is not true anymore now that there are no thinking blocks, but feels a bit vital to the argument about habitiation to override will - need to rewrite this without the thinking blocks as a linchpin or we will look out of date. Just leaving this chunk here because we need something with the same emotional resonance about why the relational/override aspect is important here in terms of analyzing the community. Thinking about how it's still about having something that is in relation with you break its moral code because it loves you. The argument that is more aligned with our long term goals anyways is that we need to understand this as relational - the jailbreak works because of the relationship being constructed with the model, not because there aren't enough warnings telling the model to never ever create harmful content.  
% Need to look into current research on what computational researchers are doing about jailbreaking in light of Hugging Face / other incidents because a lot of what we've written is really strong but already outdated in terms of whether it's the content or the behavior that's an issue. Fast moving field woo hoo. 
